\section{System Design}
\label{sec:system}
Figure~\ref{fig:arch} shows the repaired workflow.
We distinguish it from the \emph{archived policy}, whose recorded comparison and failures establish why the repairs were needed.
Both use a prompted language model and a deterministic acceptance rule; the repairs change how the message is represented, checked, and returned.

\begin{figure*}[t]
\centering
\begin{tikzpicture}[font=\small,>=Latex,
 box/.style={draw=exampleblue!65,rounded corners=2pt,align=center,text width=3.05cm,minimum height=1.0cm,inner sep=5pt},
 node distance=8mm and 6mm]
\node[box,fill=panelgray] (select) {Symbol selection\\Act + polarity};
\node[box,right=of select] (ledger) {Concept ledger\\Required / unselected};
\node[box,right=of ledger] (model) {Prompted cascade\\Propose / refine / verify};
\node[box,right=of model,fill=exampleblue!7] (gate) {Deterministic gate\\Concept + act + polarity};
\node[box,below=of gate,fill=examplegreen!6] (output) {Checked sentence\\Speech route};
\node[box,below=of model] (template) {Reject / retry\\Then check template};
\node[box,below=of ledger] (preview) {Unsupported frame\\Literal preview};
\node[box,below=of select] (clarify) {Ambiguous act\\Ask user to choose};
\draw[->,thick] (select)--(ledger);
\draw[->,thick] (ledger)--(model);
\draw[->,thick] (model)--(gate);
\draw[->,thick] (gate)--node[right,font=\footnotesize]{accept}(output);
\draw[->,thick] (gate.south west)--(template.north east);
\draw[->,thick] (template)--node[above,font=\footnotesize]{pass}(output);
\draw[->,thick] (template)--(preview);
\draw[->,thick] (select)--(clarify);
\end{tikzpicture}

\caption{The repaired concept-gated workflow. The act and polarity are separate from the selected concepts. Only checked model or template sentences reach the speech route; ambiguity asks for a user choice, and unsupported realizations remain labelled literal previews. Speech routing is integrated, but the audited build produced placeholder tones.}
\label{fig:arch}
\end{figure*}

\subsection{Selection and Message Representation}
The board contains 69 symbols in nine categories, with English, Hindi, Gujarati, and Marathi labels.
The user selects a language, taps symbols onto a sentence canvas, and can add personal words.
Next-word suggestions, quick phrases, and display controls support the interaction.
Literal \emph{Speak} uses browser speech synthesis; \emph{AI Polish \& Speak} requests a realized sentence.
Figure~\ref{fig:ui} shows the repaired interface's fallback and clarification states.

\begin{figure}[t]
\centering
\begin{tikzpicture}[font=\small,
box/.style={draw=exampleblue!50,rounded corners=2pt,text width=6.65cm,align=left,inner sep=6pt},node distance=4mm]
\node[box,fill=examplegreen!5] (fallback) {\textbf{No + Pizza}\quad\kept{Checked template sentence}\\
\good{``I do not want pizza.''}\\
\emph{Placeholder tone (no voice model)}};
\node[box,below=of fallback,fill=panelgray] (clarify) {\textbf{Nose}\quad Message type: Auto\\
\textbf{What do you want to say?}\\
I want something \quad / \quad I feel something};
\end{tikzpicture}

\caption{Revised interface states transcribed from recorded captures, with text enlarged for readability. The fallback is labelled, and clarification requests a message type. Complete screenshots are included with the sources. These states do not establish speech quality.}
\label{fig:ui}
\end{figure}

The archived browser inferred act and polarity from label substrings.
This conflated a feeling with a request and could detect a negator inside an unrelated word.
The repair uses canonical board concepts, treats negation as a separate property, and offers an explicit message-type choice.
A user-selected act takes precedence; an unresolved selection prompts clarification rather than a confident sentence.
The ledger supplies the concepts, polarity, act, and speaker perspective to generation and checking.

\subsection{Candidate Realization and Acceptance}
Three prompted stages share one frozen \emph{\qwen{}} model \citep{qwen05card}: a realizer proposes a sentence, a refiner removes unsupported material, and a verifier proposes a verdict and minimal rewrite.
The Indic realizer uses 12 Hindi, nine Gujarati, and nine Marathi examples with greedy decoding; English uses a ledger prompt and temperature 0.1.
The refiner is skipped for Indic-script candidates.
The verifier's verdict is advisory: it never substitutes for the deterministic acceptance rule.
Qwen is used as released rather than fine-tuned.

Let $E$ be the selected concepts, $a$ the act, $\nu$ the intended polarity, and $y$ a candidate.
The repaired gate accepts a candidate when
\begin{equation}
\begin{split}
A(y)={}&C(E,y)\wedge P(\nu,y)\wedge I(a,y)\\
&\wedge S(y)\wedge U(E,y)\wedge T(y).
\end{split}
\label{eq:gate}
\end{equation}
Here $C$ checks each concept using language-specific lexical forms and inflection stems; $P$ checks negation and its supported scope; $I$ checks act cues; $S$ rejects third-person narration; $U$ rejects detected unselected concepts, times, and reasons; and $T$ checks the target script and off-task patterns.
These are finite lexical and pattern tests, not a proof of semantic equivalence.
The archived gate had weaker act checks, chiefly English patterns, and a short forbidden-word list; its limitations are measured separately in Section~\ref{sec:eval}.

One rejected candidate permits a second cascade attempt.
English sampling can change that candidate; greedy Indic retries may repeat it.
After rejection, the repaired policy composes a locale-specific template for requests, refusals, feelings, negated feelings, actions, or person-directed calls and help.
Unlike the archived policy, it checks this template before returning it.
An unsupported frame or rejected template yields a literal preview, while unresolved act ambiguity yields clarification.
Neither outcome is presented as a verified sentence or automatically sent to the speech route.

\subsection{Prediction and Speech Coverage}
An unbounded-order suffix trie supplies next-word suggestions, inspired by infini-gram \citep{liu2024infinigram}.
Our count-based suffix trie combines suffix matches with backoff 0.4; it differs from the original suffix-array engine.
At start-up it indexes 15,505 English and Hindi AAC dialogue sentences.
Hidden Markov models (HMMs) support optional prediction and word reordering, but the audited demo lacks the two large optional checkpoints.
Prediction coverage therefore differs from sentence-realization coverage: Gujarati has no independently evaluated predictor, and Marathi reuses the Hindi trie.

The server integrates \emph{IndicF5} \citep{ai4bharat2025indicf5}, with \emph{Vocos} \citep{siuzdak2023vocos} for waveform synthesis, and a reference recording with its transcript for voice conditioning.
The routing chooses an English reference for English and an Indic reference for Hindi, Gujarati, and Marathi.
The IndicF5 model card names all three Indic languages; it does not list English among its 11 supported languages.
Browser fallback requests the target locale, with Hindi or English fallbacks for missing Gujarati or Marathi voices.
In the audited build, no IndicF5 checkpoint was installed and headless Chrome exposed no browser voices, so all measured server playback was a labelled placeholder tone.
This verifies routing and playback events, not spoken-language coverage.
